\section{The homogeneous four-family calculus}
\label{sec:homogeneous}

The proposed calculus retains the scalar terms and the original ordinary
meaning of the repository's Eq/$D$ language.  The change is in retained
assumptions: every context entry is a single one-variable formula
template.  No entry has a separately stored antecedent, a list of
antecedents, or a flag distinguishing facts from conditional closures.

\subsection{Language, ordinary reading, and the gate}

The term language is
\[
 t ::= x\mid 0\mid St\mid I(n,a;u.h),
\]
where $u$ is bound only in $h$.  Formulas are
\[
 A ::= (s=t)\mid D_{x,y}A.
\]
The ordinary reading is unchanged:
\[
 D_{x,y}B(x,y)\ :=\ \exists x\,
             \bigl(B(x,0)\to\forall y\,B(x,y)\bigr).
\]
Write $Q_yB$ as a metalinguistic abbreviation for
$B(0)\to\forall yB(y)$, not as a constructor of the new formula language.
A context $\Delta$ is a finite reusable list of templates $B(y)$; its
ordinary reading is the list of formulas $Q_yB$.  The distinguished
initial entry is
\[
 \rho(y)\equiv(y=y).
\]
It remains present in every context.  Its ordinary reading is a theorem.

For a fresh variable $v$, define
\[
 \kappa_v(s=t)
   :=\bigl(I(0,s;u.t)=I(v,s;u.t)\bigr),
 \qquad
 \kappa_v(D_{x,y}B):=D_{x,y}(\kappa_vB).
\]
Here $u$ is fresh and unused in $t$; all substitutions and insertions are
capture-free.  In a de Bruijn implementation, placing $t$ in the new
iterator body must lift its free indices.  The asymmetry is useful:
\[
 \kappa_0(s=t)\text{ is literally reflexive},\qquad
 \kappa_{Sz}(s=t)\ \dashv\vdash\ (s=t).
\]
The second statement uses iterator computation and identifies the
positive gate with the original equation.  Under nested
$D$ binders, the outside variable $v$ remains free; the gate does not
change the number of formula binders or an equation/compound root.

\subsection{The four rule families}

\paragraph{Use.}
For an entry $B(y)$ and a target term $t$, infer $B(t)$ from one minor
derivation of $B(0)$.  The minor may instead be omitted exactly when the
raw substituted base $B(0)$ has the form $D^j(r=r)$ with identical
leaf terms, for some $j\ge0$:
\[
 \frac{B(y)\in\Delta\qquad\Delta\vdash_H B(0)}
      {\Delta\vdash_H B(t)},
 \qquad
 \frac{B(y)\in\Delta\qquad B(0)\equiv D^j(r=r)}
      {\Delta\vdash_H B(t)}.
 \tag{Use}
\]
These are the one-premise and zero-premise cases of one instruction
family.  The omission test compares the two literal leaf terms and
does not inspect their numerical values.  It is valid at every binder
depth; an arbitrary compound theorem does not qualify.

\paragraph{Eval.}
From two equations, infer the result of one permitted computation:
\[
 \frac{\Delta\vdash_H s=t\qquad\Delta\vdash_H b=r}
      {\Delta\vdash_H r=\Phi(s,t,b)}.
 \tag{Eval}
\]
The partial operation $\Phi$ has exactly two cases:
\[
 \Phi(s,t,b)=
 \begin{cases}
 t,&b\equiv I(0,s;u.h),\\
 h[t/u],&b\equiv I(Sc,a;u.h)\ \text{and }s\equiv I(c,a;u.h).
 \end{cases}
\]
All other shapes are rejected.  In particular, no primitive root-joining
or arbitrary transport case is included.

\paragraph{Elim.}
\[
 \frac{\Delta\vdash_H D_{x,y}B(x,y)\qquad
       \Delta,B(b,y)\vdash_H A}
      {\Delta\vdash_H A}.
 \tag{Elim}
\]
The eigenparameter $b$ is fresh for the conclusion, original context and
fixed parameters.  The new retained entry is the single template
$y\mapsto B(b,y)$, with no separately stored zero instance.

\paragraph{Intro.}
Define
\[
 \eta(B,t,z)=
 \begin{cases}
 z,&B\text{ is a literal equation and }t\equiv0,\\
 0,&\text{otherwise}.
 \end{cases}
\]
Then
\[
 \frac{\Delta,\kappa_v(B(t,\eta(B,t,z)))\vdash_H B(t,Sz)}
      {\Delta\vdash_H D_{x,y}B(x,y)}.
 \tag{Intro}
\]
The eigenparameter $z$ is fresh for $\Delta,t$ and all fixed parameters
of $B$.  The displayed new entry is a template in its fresh argument $v$;
$z$ is a free parameter there.  The witness-directed choice of cases or
progress is exactly the repository's established witness rule.

There are no further primitive rules or nonlogical axioms.  In
particular, there is no Hyp rule and no Refl rule.  The designated
contradiction is $0=S0$.

\subsection{Recovering computation and numerical equality}

\paragraph{Reflexivity.}
Use the fixed entry $\rho(y)=(y=y)$ at target $t$.  Its base is literally
$0=0$, so no minor is required and the conclusion is $t=t$.


\paragraph{Reflexive prefixes.}
Every well-scoped $D^j(r=r)$ is derivable in every context.  Prove this
by induction on $j$, independently of the gate retrieval theorem.  The
case $j=0$ is the preceding derived reflexivity.  For $j>0$, apply Intro
with witness $1=S0$, which forces case mode.  Ignore its temporary gated
assumption.  The required successor target has $j-1$ binders, and its
leaf terms remain identical because both undergo the same substitutions.
The induction hypothesis supplies that target.  Consequently every
omitted Use minor can be expanded into an actual derivation using only
the fixed initial matrix and the displayed rules.

\paragraph{The two computations.}
Zero computation is Eval on the reflexivities $a=a$ and
$I(0,a;u.h)=I(0,a;u.h)$.  Successor computation is Eval on the
reflexivities of $I(c,a;u.h)$ and $I(Sc,a;u.h)$.

\paragraph{Root joining.}
From $s=t$ and $s=r$, put $Z=I(0,s;u.0)$.  Eval on $s=r$ and $Z=Z$
gives $Z=r$; Eval on $s=t$ and $Z=r$ gives $r=t$.  Symmetry is this
derived joining rule with one premise reflexive; transitivity follows by
one symmetry and one joining.

\paragraph{Transport in a term template.}
Suppose $s=t$ and $h(s)=r$.  Set
\[
 Z=I(0,s;u.h(u)),\qquad P=I(S0,s;u.h(u)).
\]
Zero-case Eval yields $Z=s$ and $Z=t$.  Successor-case Eval on $Z=s$
and $P=P$ gives $P=h(s)$, hence $P=r$ by the preceding derived equality
rules.  Successor-case Eval on $Z=t$ and $P=r$ now yields $r=h(t)$.
This is the native kernel's crossed one-template transport rule.
Consequently all numerical equality reasoning used below is derived
inside the four-family calculus.

\subsection{The gate retrieval theorem}
\label{subsec:gate-retrieval}

For every formula $A$ and every context containing $\rho$:
\begin{enumerate}
\item $\Delta\vdash_H\kappa_0A$.
\item A derivation of $\kappa_1A$ can be converted effectively to a
      derivation of $A$, and conversely.
\item If $\kappa_vA$ is a retained context entry, then
      $\Delta\vdash_H A$.
\end{enumerate}
The transformations are uniform under enlarging the context and under
renaming eigenparameters.

\paragraph{Proof of assertion 1.}
The formula $\kappa_0A$ has the same $D$ prefix as $A$ and literally
identical iterator terms at its leaf.  It is therefore an instance of
the independently established reflexive-prefix theorem.  No retrieval
lemma is used in this step.

\paragraph{Simultaneous proof of assertions 2 and 3.}
At an equation, assertion 2 is computation and equational transport.  Suppose assertion 3 has been proved for all formulas with
fewer binders than $A=D_{x,y}B$.

To convert $\kappa_1A=D_{x,y}\kappa_1B$ to $A$, eliminate the source
at a fresh $b$, retaining the matrix $\kappa_1B(b,y)$.  Introduce the
target with witness $b$.  Because $b$ is a variable rather than the
literal zero, this is a case introduction.  Its local entry is
$\kappa_v(B(b,0))$.  Assertion 3 at the smaller formula depth supplies
$B(b,0)$.  Assertion 2 at that smaller depth converts it to
$\kappa_1B(b,0)$.  Use the retained source matrix at target $Sz$ and
convert its result $\kappa_1B(b,Sz)$ back to $B(b,Sz)$.  Close Intro
and Elim.

The reverse conversion starts by eliminating $D_{x,y}B$ and introducing
$D_{x,y}\kappa_1B$ with the same fresh witness $b$.  Retrieve its local
zero assumption by assertion 3 at the smaller depth, convert to the
ungated zero instance, Use the retained $B(b,y)$ matrix at $Sz$, and
convert back to the gated successor instance.  Thus assertion 2 at the
current depth uses assertion 3 only at a strictly smaller depth.

Finally, prove assertion 3 at the current depth.  Use the retained entry
$\kappa_vA$ at target $1$.  Its zero instance is leaf-reflexive at every
formula depth, so the minor is omitted by the explicitly stated Use
rule.  The result is $\kappa_1A$, which assertion 2 converts to $A$.
This completes a well-founded simultaneous induction.  If desired, the
omitted minor can be expanded by assertion 1; no unproved compound
assumption has been discarded.

\subsection{Soundness in the original ordinary realization}

Read context entries as their conditional closures in $T_{\rm eq}$.
Formula induction in ordinary intuitionistic logic with equality proves
\[
 \kappa_0A,\qquad \kappa_{Sz}A\leftrightarrow A.
\]
At atoms the latter uses the two computations, and at a $D$ binder it
uses the ordinary extensionality of the displayed existential conditional
definition.  Numerical coverage therefore proves
\[
 A\ \longrightarrow\ Q_v\kappa_vA.
 \tag{Gate introduction}
\]
Indeed, under $A$, a numerical argument is zero or a successor; the
first gated instance is a theorem, and the second is equivalent to $A$.

Induct now on a four-family derivation.  Eval is a valid equational
consequence.  Use is conditional-universal specialization; when its
minor is omitted, the reflexive-prefix theorem supplies exactly that
missing minor.  Elim is ordinary existential elimination.  For Intro, its
ordinary temporary fact $B(t,\eta)$ supplies the retained encoded
assumption by Gate introduction.  When $\eta=0$, the successor premise
and numerical coverage establish the conditional universal conclusion.
When $\eta=z$, the matrix is an equation, and equation induction with
retained side assumptions establishes the same conditional conclusion.
Existential introduction at $t$ completes the argument.  The initial
$\rho$ closure is an equality theorem.  Hence the calculus is sound
under the unchanged ordinary reading.

\subsection{Completeness and the interpretation composites}

Let $W$ be the pinned seven-rule witness calculus \texttt{WPrf}.  The
repository proves \texttt{wprf\_iff\_nprf}; its introduction uses the
same witness test as above.  Translate its formulas by fixing equations
and $D$ and putting
\[
 (Q_yB)^\circ=D_{r,y}B^\circ
\]
with $r$ fresh and unused.  The translation preserves root class, fixed
parameters and capture-free numerical substitution.

The precise derivation induction is: if $\Gamma\vdash_W A$, and every
$C\in\Gamma$ has a proof of $C^\circ$ from a homogeneous context
$\Delta$, then $\Delta\vdash_H A^\circ$.  Hyp uses the supplied proof.
The computations and transport use the equational derivations above.
For $Q$ elimination, eliminate the translated dummy-witness $D$ and Use
its retained matrix at the original target.  For Intro, the gate
retrieval theorem supplies precisely the translated local fact from the
new entry $\kappa_v(B^\circ(t,\eta))$, so the induction hypothesis
applies to the source premise.

For source $D$ elimination, the newly retained matrix is
$B^\circ(b,y)$.  It supplies the translated source assumption
$(Q_yB(b,y))^\circ=D_{r,y}B^\circ(b,y)$ as follows: introduce this
dummy-witness binder at witness $1$, ensuring case mode, retrieve its
local zero fact through the gate theorem, and Use the retained matrix at
$Sz$.  Thus the derivation induction also applies to the source
elimination premise.  This strengthened context statement avoids any
unproved primitive cut or clause-construction rule.

For an explicit source context, take $\Delta$ to consist of $\rho$ and
the entries $\kappa_v(C^\circ)$ for $C\in\Gamma$; gate retrieval
provides all needed hypothesis proofs.  Hence every source derivation
has an effective four-family derivation.

Conversely, the soundness proof translates every new rule to the original
ordinary scalar theory.  The pinned native completeness theorem returns
it to the native calculus.  For context bookkeeping, the reading of a
matrix entry is exactly its native $Q$ formula, and
$Q_v\kappa_vA\leftrightarrow A$ follows from Gate introduction and
specialization at $1$.  Thus the encoded-context composites are also
provably equivalent, not merely equiconsistent.

On formulas, the $H$-to-native translation is the inclusion of the
Eq/$D$ language.  The other translation only replaces $Q$ by $D$ with
an unused witness.  The native theory proves this replacement
equivalent in both directions, so formula induction gives both
round trips.  The $H$-side formula composite is literally the identity.
Both term maps and both numerical domains are the identity; numerical
equality remains actual equality.  Composing with the repository's
native-to-iPRA bi-interpretation gives the requested bi-interpretability,
both numerical composite graphs $x=y$, and equiconsistency.  The
designated contradiction is fixed.  The arithmetic bridge and its
existing proof-assistant verification are inherited results; the new
gate and context theorem is a written proof unless separately formalized.

\subsection{A flat formula representation}

The Eq/$D$ grammar is unary above its atomic leaf.  Every formula is
uniquely of the form
\[
 D_{x_1,y_1}\cdots D_{x_k,y_k}(s=t).
\]
It may therefore be represented by $(k,s,t)$ instead of a recursively
nested formula tree, with the terms scoped under the $2k$ displayed
bound variables.  A gate leaves $k$ unchanged and transforms the single
leaf equation; a binder adds one to $k$.  Under a relative-index
representation, the gate's external variable must be placed outside all
$2k$ formula-bound variables, and placing the old right-hand term under the iterator binder
requires lifting its exterior positive indices. In the mixed convention
of the delivered checker, negative free-parameter levels remain unchanged.  This is a representation theorem with no change
of language or interpretation.  Checker measurements must still include
the integer and scope validation, capture-free substitutions and gate
construction.

\paragraph{Scope of the simplicity claim.}
The mathematical deletion is the separate antecedent field of every
retained clause.  Fact assumptions and conditional closures now use one
homogeneous context grammar.  The omitted-minor reflexivity test and the
gate are explicit parts of the trusted checker.  Four instruction
families do not beat the repository's existing four-family count by
themselves; any software-size improvement must be the result of a
matched, complete, measured implementation.


\paragraph{Implementation evidence.}
The separate homogeneous implementation was checked on 400 generated
formulas for its literally reflexive zero gate and on 800 positive-control
instances against environment-based numerical evaluation. Six closed
certificates passed, including equation induction whose local hypothesis
is recovered from its gated context entry, followed by specialization.
These finite checks are reproduced by \texttt{code/audit\_homogeneous.py};
they are separate from the primary checker's subprocess comparison.
